% !TEX root = dkls23.tex
% title: DKLS23 论文精读


\subsubsection{pp3 sec1.1“A Brief History of Threshold ECDSA”}\label{pp3-sec1.1-a-brief-history-of-threshold-ecdsa}

每种历史方案 = 改写 ECDSA 签名方程的方式 + 针对改写设计的算法 +
针对改写设计的防作弊检查

\subsubsection{pp5 sec1.1“bespoke”}\label{pp5-sec1.1-bespoke}

定制，专用

\subsubsection{pp6 sec1.1“Generic Approaches”}\label{pp6-sec1.1-generic-approaches}

不为 ECDSA 设计专门的正确性和防作弊协议，而是把算法丢进 MPC 电路，如
SPDZ。电路里的每一步都需要做通用的 + 与具体问题无关的认证，
比专用协议开销大得多。

相关:“a factor of two in terms of bandwidth and computation”

\subsubsection{pp7 sec1.2“Rewriting ECDSA”}\label{pp7-sec1.2-rewriting-ecdsa}

\begin{verbatim}
R := r * G;
a := Hash(msg);
b := R.x;
w := (a + sk*b) * phi;
u := r * phi
s := w / u;
return R, s;
\end{verbatim}

\subsubsection{\texorpdfstring{pp7 sec1.2“secret shares $w$”}{pp7 sec1.2“secret sharesw”}}\label{pp7-sec1.2-secret-sharesw}

原句：The computation of secret shares of $w$ can be performed locally
by the parties given shares $\phi$ and $v = \mathtt{sk}\cdot\phi$。

这句话的意思是，
在各方已经摇出 $\phi$ 的加法分片以及算出 $v$ 的加法分片后，
各方就可以本地合成 $w$ 的加法分片。原理如下式：
\begin{equation}
\label{eq:dkls23:1}
\begin{aligned}
w &= (a+\mathtt{sk}\cdot b)\phi \\
&= a\phi + b\cdot \mathtt{sk}\cdot \phi \\
&= a\phi + bv.
\end{aligned}
\end{equation}

\subsubsection{\texorpdfstring{pp7 sec1.2“$v$ and $u$ are computed ideally”}{pp7 sec1.2“vanduare computed ideally”}}\label{pp7-sec1.2-vanduare-computed-ideally}

原句：Assuming that shares of the two products $v$ and $u$ are computed
ideally \ldots{}

这句话的的意思是，计算 $u$ 和 $v$ 的加法分片需要用到 MtA 协议。
其中，$u = r\phi$，$v = \mathtt{sk}\cdot\phi$。

\subsubsection{pp7，8 sec1.2“a few avenues to cheat”}\label{pp78-sec1.2-a-few-avenues-to-cheat}

假设 MtA 本身安全，敌手还能靠什么歪门邪道搞破坏? 只有以下四条路：

\begin{enumerate}
\def\labelenumi{(\arabic{enumi})}

\item
  有偏的 nonce。如果某个敌手能看到别人的 $r_i$ 分片，
  适应性地调整自己的 $r_i$，就能操纵最终 $R$ 的分布。
  进而通过格攻击等手段泄露私钥。
\end{enumerate}

防御：commit-then-reveal。各方先各自采样 $r_i$，
为相应的 $R_i$ 发布承诺。等所有人都承诺完了再一起解承诺。
只要 $r$ 在承诺阶段结束前是信息论隐藏的，
敌手就不可能让自己的 $R_j$ 依赖诚实方的 $R_i$。

（2）$v$ 和 $u$ 使用不一致的 $\phi$。如果某个敌手使用不一致的 $\phi_j$，
那么会导致验签失败。

防御：不设置检查，而是从结构上消灭使用两个 $\phi$ 的机会。详见本文
Vector OLE 相关内容。

\begin{enumerate}
\def\labelenumi{(\arabic{enumi})}
\setcounter{enumi}{2}

\item
  使用不一致的 $\mathtt{sk}_i$ 或 $r_i$。
  如果某个敌手使用的 $\mathtt{sk}_j$ 或 $r_j$ 与公开承诺不一致，
  那么验签失败。
\end{enumerate}

防御：利用 $\phi_i$ 作为顺带的 MAC key。具体来说，MtA
算出的 $r_j \cdot \phi_i$ 顺带就是“以 $\phi_i$ 为密钥，给 $r_j$ 打的
MAC”。验证挪到曲线上做。

\begin{enumerate}
\def\labelenumi{(\arabic{enumi})}
\setcounter{enumi}{3}

\item
  直接篡改 $v$，$u$。后果是验签失败。
\end{enumerate}

防御：因为协议跑到这里已经快要结束，所以交给验签去拦截。

追问：途径 2 和 3 的后果也是验签失败，为什么不能像 4
一样甩给最后一步验签?

回答：论文要防的远不止最终签名对不对，还包括：整个协议执行过程中，
没有任何一步会把诚实方的秘密泄露出去（UC/模拟安全）。

\begin{verbatim}
作弊发生时, 协议必须 abort; 否则诚实方会因不知情而把更多信息交给作弊方.
\end{verbatim}

途径 4 能够甩给验签去拦截，也是因为诚实方已经没有信息可以交出去。

\subsubsection{什么是 MAC}

MAC = Message Authentication Code。假设 Alice 和 Bob
共享一个只有他俩知道的密钥 $k$。Alice 想发一条消息 $m$，又怕被篡改，
于是她算一个标签 $t := \mathrm{Hash}(k, m)$，把 $m, t$ 一起发给 Bob。
然后 Bob 重放 $t$，看跟收到的 $t$ 是否一致。

\subsubsection{pp11 sec2.2“malicious PPT adversary”}\label{pp11-sec2.2-malicious-ppt-adversary}

原句：We consider a malicious PPT adversary who can statically corrupt
up to $t - 1$ parties。

我的疑问：corrupted parties 腐化或诚实取决于他们的自由意志。
他们或许可以把责任推给一个 adversary，但 adversary 为什么是 PPT 的?
这让我觉得无厘头，就好像说一位肇事司机是 PPT。

PPT adversary $\mathcal{A}$ 是“所有被腐化方合谋能干的坏事”
的数学代号。

原句的每一个定（状）语都有明确的意图：* malicious。坏人能不能任意偏离协议，
乱发消息? 能。* PPT。坏人算力有多强? Probabilistic polynomial-time。*
statically。坏人什么时候决定拉拢谁? 协议开始前定好，中途既不反悔，
也不追加。* up to $t-1$。坏人最多能拉拢多少人? 不超过 $t-1$ 人。*
单一的 $\mathcal{A}$。坏人之间能不能互相通气? 能，
而且安全证明要覆盖最大限度的协同作案。

类比肇事司机的例子。司机的意志/策略空间是无限自由的（定语 malicious）。
但他的车最快只能跑200码（定语 PPT）。虽然司机的策略空间是无限自由的，
但为了方便分析，我们追加假设司机一旦动手就不收手（状语 statically）。

\subsubsection{pp12 sec3“This distinction”所在的一整段}

我们的（功能定义）把 ECDSA 算法当作黑盒使用，不会提前泄露 R，
并且在形式上把 abort（阻断此后与该功能的一切交互）同 failed signature
（不会阻断交互）区分开来。这个区分在门限（签名）功能定义里很重要，
因为不应该出现这种情况：仅仅因为某个腐败方参与了一次签名，
以后不带他也不能签名。

以上说的是 DKLs19 的缺陷。

\subsubsection{什么是理想功能}

“这个协议安全吗?”是一句很模糊的话，安全指什么? 不泄密? 不能被伪造?
不能被拒绝服务? 这个清单用自然语言是列不完的，
几乎必然漏掉某个没想到的边角案例。

UC 框架换了个思路。不列清单，而是虚构一个绝对靠谱的第三方来完成这个任务。
这个虚构的第三方就叫理想功能（ideal functionality）。

理想功能遵循通用模板：* 参数。比如群 G，参与方数量 n，门限 t。*
跟诚实的参与方打交道.“当从某方 Pi 收到某消息时，就做某事”* 跟敌手 S
（Simulator）打交道。提供给 S 一份有限的命令清单。这是 S
做坏事的能力边界。

我能理解，对于诚实参与方来说，理想功能很像接口 / 虚函数。但我不理解，
凭什么敌手 S 只能运用有限的接口?

实际上，理想功能并没有限制现实世界的坏蛋 $\mathcal{A}$，
而是在做一个数学断言：

不管 $\mathcal{A}$ 在现实里有多离谱，捣乱方法多么天马行空，
它给诚实方造成的后果总能被归类为几个有限的命令，不多不少。

UC 证明并不是彰显协议固若金汤的证书。它实际上把“协议安全吗”
这个模糊的问题，严格地拆分为两块（三块）边界分明的责任：1。
协议逻辑安全吗? 由论文的数学证明来保证。一旦证完，几乎不可能是事故来源。
2。密码学假设安全吗? 由更基础的理论来保证。有可能被攻破，但理论门槛极高。
3。模型实现安全吗? 由实现和部署的质量来保证。有可能留下漏洞，
如随机数复用，旁路信道等。

UC 框架还需要 Hybrid 证明，它的职责是证明“真实协议+$\mathcal{A}$”
和“理想功能+$\mathcal{S}$”无法区分。具体来说，

\begin{enumerate}
\def\labelenumi{(\arabic{enumi})}

\item
  谁跟谁无法区分?
\end{enumerate}

环境 $\mathcal{Z}$ 与（真实协议 +$\mathcal{A}$）的交互输出，
与 $\mathcal{Z}$ 与（理想功能 +$\mathcal{S}$）的交互输出，
这两个交互输出无法区分。

$\mathcal{Z}$ 没有具体的内在结构，它就是“任意一台交互式图灵机”。
研究者如果追求计算安全，那就限定它的算力为 PPT。如果追求统计安全，
则可以无界。

既然 $\mathcal{Z}$ 是“任意的”，那么我们只能规定它的接口清单：1。
给诚实方写输入：如选哪条消息签名，选哪个门限，何时调 Setup（keygen），
何时调 Sign。而且可以在交互途中自适应地选。2。从诚实方读输出：如公钥，
签名，abort 信号。3。与敌手自由通信。$\mathcal{Z}$ 通过

\begin{enumerate}
\def\labelenumi{(\arabic{enumi})}
\setcounter{enumi}{1}
\item
  为什么要无法区分? TODO
\item
  如何衡量无法区分? Hybrid
  证明的思路是在这两个交互输出之间安插若干里程碑。
  上一个里程碑相应的协议修改一处就成为当前里程碑。根据三角不等式，
  两端的统计距离不大于相邻里程碑统计距离之和。
\end{enumerate}

\subsubsection{pp12 Functionality 3.1}\label{pp12-functionality-3.1}

“If any party is corrupt \ldots{} may be attempted”
